\documentclass[11pt, a4paper]{article}
\usepackage[utf8]{inputenc}
\usepackage[margin=0.75in]{geometry}
\usepackage{amsmath, amssymb}
\usepackage{enumitem}
\usepackage{multicol}

\title{\textbf{NEET Biology: Respiratory Quotient (RQ) — High-Yield Difficult Questions}}
\author{Target: NEET UG}
\date{}

\begin{document}

\maketitle

\section*{Questions}

\begin{enumerate}[label=\textbf{\arabic*.}]

\item Which of the following equations accurately defines the Respiratory Quotient ($\text{RQ}$) of an organism or tissue?
\begin{enumerate}[label=(\Alph*)]
    \item $\text{RQ} = \dfrac{\text{Volume of } \text{O}_2 \text{ consumed}}{\text{Volume of } \text{CO}_2 \text{ evolved}}$
    \item $\text{RQ} = \dfrac{\text{Volume of } \text{CO}_2 \text{ evolved}}{\text{Volume of } \text{O}_2 \text{ consumed}}$
    \item $\text{RQ} = \dfrac{\text{Moles of } \text{ATP} \text{ synthesized}}{\text{Moles of } \text{O}_2 \text{ consumed}}$
    \item $\text{RQ} = \dfrac{\text{Mass of substrate oxidized}}{\text{Volume of } \text{CO}_2 \text{ evolved}}$
\end{enumerate}

\item What is the exact Respiratory Quotient ($\text{RQ}$) for the complete aerobic oxidation of the triglyceride **tripalmitin** ($\text{C}_{51}\text{H}_{98}\text{O}_6$), as given in NCERT Class 11 Biology?
\begin{enumerate}[label=(\Alph*)]
    \item $1.0$
    \item $0.9$
    \item $0.7$
    \item $1.33$
\end{enumerate}

\item Consider the balanced chemical equation for the complete aerobic respiration of **tripalmitin**:
$$2(\text{C}_{51}\text{H}_{98}\text{O}_6) + x\text{O}_2 \longrightarrow y\text{CO}_2 + z\text{H}_2\text{O} + \text{Energy}$$
Identify the exact stoichiometric values of $x$, $y$, and $z$:
\begin{enumerate}[label=(\Alph*)]
    \item $x = 145$, $y = 102$, $z = 98$
    \item $x = 102$, $y = 145$, $z = 98$
    \item $x = 72.5$, $y = 51$, $z = 49$
    \item $x = 145$, $y = 98$, $z = 102$
\end{enumerate}

\item \textbf{Assertion (A):} The Respiratory Quotient ($\text{RQ}$) of fats and fatty acids is always less than unity ($\approx 0.7$). \\
\textbf{Reason (R):} Fats contain a much lower proportion of oxygen relative to carbon and hydrogen compared to carbohydrates, and therefore require a larger volume of external $\text{O}_2$ for complete oxidation to $\text{CO}_2$ and $\text{H}_2\text{O}$.
\begin{enumerate}[label=(\Alph*)]
    \item Both (A) and (R) are true and (R) is the correct explanation of (A).
    \item Both (A) and (R) are true but (R) is not the correct explanation of (A).
    \item (A) is true but (R) is false.
    \item Both (A) and (R) are false.
\end{enumerate}

\item Calculate the theoretical Respiratory Quotient ($\text{RQ}$) for the complete aerobic oxidation of **malic acid** ($\text{C}_4\text{H}_6\text{O}_5$):
\begin{enumerate}[label=(\Alph*)]
    \item $0.70$
    \item $0.90$
    \item $1.00$
    \item $1.33$
\end{enumerate}

\item What is the Respiratory Quotient ($\text{RQ}$) of **oxalic acid** ($\text{C}_2\text{H}_2\text{O}_4$) when completely oxidized aerobically?
\begin{enumerate}[label=(\Alph*)]
    \item $1.33$
    \item $2.00$
    \item $4.00$
    \item $\infty$
\end{enumerate}

\item Match Column-I (Respiratory Substrate / Condition) with Column-II (Typical $\text{RQ}$ Value):
\begin{center}
\begin{tabular}{|l|l|}
\hline
\textbf{Column-I} & \textbf{Column-II} \\
\hline
P. Germinating wheat seeds (starchy) & 1. $0.7$ \\
Q. Germinating castor or groundnut seeds (fatty) & 2. $0.9$ \\
R. Germinating pea or bean seeds (proteinaceous) & 3. $1.0$ \\
S. *Opuntia* (CAM plant) at night & 4. $0.0$ \\
\hline
\end{tabular}
\end{center}
\begin{enumerate}[label=(\Alph*)]
    \item P-3, Q-1, R-2, S-4
    \item P-3, Q-2, R-1, S-4
    \item P-1, Q-3, R-2, S-4
    \item P-4, Q-1, R-2, S-3
\end{enumerate}

\item Why is the Respiratory Quotient ($\text{RQ}$) of succulent CAM plants (such as *Opuntia* and *Bryophyllum*) equal to **zero** during the night?
\begin{enumerate}[label=(\Alph*)]
    \item Neither $\text{O}_2$ is consumed nor $\text{CO}_2$ is evolved at night.
    \item Carbohydrates are incompletely oxidized into organic acids (like malic acid) using $\text{O}_2$, while all $\text{CO}_2$ is internally fixed by PEP carboxylase so no $\text{CO}_2$ is released externally.
    \item Stomata remain closed at night, trapping $\text{O}_2$ inside the mesophyll cells.
    \item CAM plants perform alcoholic fermentation exclusively at night.
\end{enumerate}

\item In *Opuntia* (a succulent CAM plant), what happens to the Respiratory Quotient ($\text{RQ}$) during the **daytime** when the organic acids accumulated overnight are decarboxylated and respired?
\begin{enumerate}[label=(\Alph*)]
    \item $\text{RQ}$ remains $0$
    \item $\text{RQ}$ becomes less than $1$ ($\approx 0.7$)
    \item $\text{RQ}$ becomes greater than $1$ ($> 1$)
    \item $\text{RQ}$ becomes equal to $\infty$
\end{enumerate}

\item What is the Respiratory Quotient ($\text{RQ}$) during **alcoholic fermentation** of glucose by *Saccharomyces cerevisiae* under strictly anaerobic conditions?
\begin{enumerate}[label=(\Alph*)]
    \item $0$
    \item $1$
    \item $2$
    \item $\infty$ (Infinity)
\end{enumerate}

\item Why is the concept of Respiratory Quotient ($\text{RQ}$) undefined or indeterminate ($0/0$) during **pure lactic acid fermentation** (such as in *Lactobacillus* or strenuously contracting human muscle)?
\begin{enumerate}[label=(\Alph*)]
    \item Because both $\text{CO}_2$ evolved and $\text{O}_2$ consumed are zero.
    \item Because $\text{O}_2$ consumed is infinite while $\text{CO}_2$ evolved is zero.
    \item Because lactic acid is immediately converted into ethanol.
    \item Because $\text{RQ}$ applies only to prokaryotic organisms.
\end{enumerate}

\item \textbf{Statement I:} In living organisms, respiratory substrates are often more than one; pure proteins or fats are never used as respiratory substrates. \\
\textbf{Statement II:} The Respiratory Quotient depends upon the type of respiratory substrate used during respiration.
\begin{enumerate}[label=(\Alph*)]
    \item Both Statement I and Statement II are correct.
    \item Both Statement I and Statement II are incorrect.
    \item Statement I is correct, but Statement II is incorrect.
    \item Statement I is incorrect, but Statement II is correct.
\end{enumerate}

\item Calculate the Respiratory Quotient ($\text{RQ}$) for the complete aerobic oxidation of **oleic acid** ($\text{C}_{18}\text{H}_{34}\text{O}_2$):
\begin{enumerate}[label=(\Alph*)]
    \item $0.706$
    \item $0.850$
    \item $1.000$
    \item $1.250$
\end{enumerate}

\item During the germination of fatty seeds (such as castor or sunflower), the $\text{RQ}$ is initially around $0.7$, then drops even lower (to $\sim 0.4\text{--}0.5$) during the peak conversion of stored triglycerides into sucrose via the **glyoxylate cycle**, and finally rises to $1.0$. Why does the $\text{RQ}$ drop below $0.7$ during active fat-to-carbohydrate conversion?
\begin{enumerate}[label=(\Alph*)]
    \item Because $\text{CO}_2$ is absorbed from the atmosphere by RuBisCO in the dark.
    \item Because converting oxygen-poor fatty acids into oxygen-rich sugars requires substantial uptake of external $\text{O}_2$ without releasing corresponding amounts of $\text{CO}_2$.
    \item Because germinating castor seeds switch to lactic acid fermentation.
    \item Because glyoxysomes consume $\text{CO}_2$ to synthesize amino acids.
\end{enumerate}

\item A germinating seed lot is placed in a respirometer and found to consume $120\text{ mL}$ of $\text{O}_2$ while evolving $84\text{ mL}$ of $\text{CO}_2$ over $2\text{ hours}$ at standard temperature and pressure. Which type of storage reserve is predominantly being oxidized by these seeds?
\begin{enumerate}[label=(\Alph*)]
    \item Starch (Carbohydrate)
    \item Proteins
    \item Fats / Oils
    \item Organic acids
\end{enumerate}

\item In **Ganong's respirometer**, if a student forgets to place potassium hydroxide ($\text{KOH}$) pellets in the small reservoir/bulb while measuring the respiration of germinating wheat seeds ($\text{RQ} = 1$), what will be observed in the graduated tube?
\begin{enumerate}[label=(\Alph*)]
    \item The level of saline/mercury will rise rapidly to the top.
    \item There will be **no net change** in the level of the liquid column because the volume of $\text{CO}_2$ evolved equals the volume of $\text{O}_2$ absorbed.
    \item The liquid column will fall rapidly due to excess gas pressure.
    \item All wheat seeds will immediately undergo alcoholic fermentation.
\end{enumerate}

\item If **Ganong's respirometer** is set up with germinating **castor seeds** ($\text{RQ} = 0.7$) **without** $\text{KOH}$ pellets in the apparatus, how will the level of the liquid column in the graduated tube behave?
\begin{enumerate}[label=(\Alph*)]
    \item It will remain completely stationary.
    \item It will **rise** slightly because the volume of $\text{O}_2$ absorbed is greater than the volume of $\text{CO}_2$ evolved, creating a partial vacuum.
    \item It will **fall** because more $\text{CO}_2$ is produced than $\text{O}_2$ consumed.
    \item It will oscillate up and down periodically.
\end{enumerate}

\item If **Ganong's respirometer** is set up with germinating **organic acid-rich tissues** ($\text{RQ} > 1$, e.g., $\text{RQ} = 1.33$) **without** $\text{KOH}$ pellets, how will the liquid column in the graduated tube move?
\begin{enumerate}[label=(\Alph*)]
    \item It will rise due to net gas consumption.
    \item It will **fall (move downward)** because the volume of $\text{CO}_2$ evolved exceeds the volume of $\text{O}_2$ absorbed, increasing internal gas pressure.
    \item It will show zero movement.
    \item It will freeze due to endothermic decarboxylation.
\end{enumerate}

\item Why is the Respiratory Quotient ($\text{RQ}$) of **proteins** approximately **$0.9$** (ranging between $0.8$ and $0.9$) rather than $1.0$ or $0.7$?
\begin{enumerate}[label=(\Alph*)]
    \item Because amino acids contain a nitrogen atom that is excreted as $\text{N}_2$ gas.
    \item Because proteins have an intermediate oxygen-to-carbon ratio compared to carbohydrates and fats, and part of the oxygen/carbon stoichiometry is affected by deamination and excretion of nitrogenous wastes (urea/uric acid).
    \item Because proteins cannot enter the Krebs cycle under any condition.
    \item Because every amino acid contains exactly $6$ carbons and $6$ oxygens.
\end{enumerate}

\item Consider a general organic compound with molecular formula $\text{C}_x\text{H}_y\text{O}_z$ undergoing complete aerobic combustion to $\text{CO}_2$ and $\text{H}_2\text{O}$. Which general formula correctly expresses its Respiratory Quotient ($\text{RQ}$)?
\begin{enumerate}[label=(\Alph*)]
    \item $\text{RQ} = \dfrac{x}{x + \frac{y}{4} - \frac{z}{2}}$
    \item $\text{RQ} = \dfrac{x + \frac{y}{4} - \frac{z}{2}}{x}$
    \item $\text{RQ} = \dfrac{y}{x + z}$
    \item $\text{RQ} = \dfrac{z}{x + \frac{y}{2}}$
\end{enumerate}

\item Using the general stoichiometric formula, calculate the Respiratory Quotient ($\text{RQ}$) of **citric acid** ($\text{C}_6\text{H}_8\text{O}_7$):
\begin{enumerate}[label=(\Alph*)]
    \item $0.75$
    \item $1.00$
    \item $1.33$
    \item $1.50$
\end{enumerate}

\item What is the Respiratory Quotient ($\text{RQ}$) of **tartaric acid** ($\text{C}_4\text{H}_6\text{O}_6$) upon complete aerobic oxidation?
\begin{enumerate}[label=(\Alph*)]
    \item $1.00$
    \item $1.33$
    \item $1.60$
    \item $2.00$
\end{enumerate}

\item Suppose an organism is oxidizing a mixture consisting of **$50\%$ carbohydrate ($\text{RQ} = 1.0$)** and **$50\%$ fat ($\text{RQ} = 0.7$)** by oxygen consumption. What will be the measured mixed Respiratory Quotient ($\text{RQ}$) of this organism?
\begin{enumerate}[label=(\Alph*)]
    \item $1.70$
    \item $0.85$
    \item $0.70$
    \item $0.35$
\end{enumerate}

\item During **starvation** in humans, the whole-body Respiratory Quotient ($\text{RQ}$) changes sequentially as different endogenous fuel reserves are depleted. Identify the correct chronological sequence of $\text{RQ}$ values from the fed state through prolonged starvation:
\begin{enumerate}[label=(\Alph*)]
    \item $0.7 \rightarrow 0.85 \rightarrow 1.0$
    \item $1.0 \text{ (carbohydrate)} \rightarrow 0.7 \text{ (adipose fat)} \rightarrow \sim 0.8\text{--}0.9 \text{ (tissue protein breakdown in terminal starvation)}$
    \item $1.33 \rightarrow 1.0 \rightarrow 0.0$
    \item $0.0 \rightarrow 0.7 \rightarrow 1.0$
\end{enumerate}

\item When carbohydrates are converted into fats (lipogenesis) in fattening animals (e.g., force-fed geese or pigs) or ripening oilseeds, what happens to the apparent Respiratory Quotient ($\text{RQ}$)?
\begin{enumerate}[label=(\Alph*)]
    \item $\text{RQ}$ becomes less than $0.5$
    \item $\text{RQ}$ becomes **greater than $1.0$** (often $1.1\text{--}1.3$) because oxygen-rich glucose is converted into oxygen-poor fatty acids, releasing extra $\text{CO}_2$ (via pyruvate dehydrogenase) without proportional $\text{O}_2$ uptake.
    \item $\text{RQ}$ drops to $0.0$
    \item $\text{RQ}$ remains strictly $0.7$
\end{enumerate}

\item \textbf{Assertion (A):} In germinating pea seeds, the $\text{RQ}$ measured during the very early stages of imbibition (before the testa ruptures) is often significantly greater than $1.0$. \\
\textbf{Reason (R):} The intact, water-imbibed seed coat restricts oxygen diffusion into the embryo, causing a combination of aerobic respiration and anaerobic alcoholic fermentation (which evolves $\text{CO}_2$ without consuming $\text{O}_2$).
\begin{enumerate}[label=(\Alph*)]
    \item Both (A) and (R) are true and (R) is the correct explanation of (A).
    \item Both (A) and (R) are true but (R) is not the correct explanation of (A).
    \item (A) is true but (R) is false.
    \item Both (A) and (R) are false.
\end{enumerate}

\item Which of the following substrates has a Respiratory Quotient ($\text{RQ}$) equal to **$1.0$**?
\begin{enumerate}[label=(\Alph*)]
    \item Sucrose ($\text{C}_{12}\text{H}_{22}\text{O}_{11}$)
    \item Palmitic acid ($\text{C}_{16}\text{H}_{32}\text{O}_2$)
    \item Malic acid ($\text{C}_4\text{H}_6\text{O}_5$)
    \item Glycine ($\text{C}_2\text{H}_5\text{NO}_2$)
\end{enumerate}

\item Calculate the Respiratory Quotient ($\text{RQ}$) for the complete aerobic oxidation of **stearic acid** ($\text{C}_{18}\text{H}_{36}\text{O}_2$):
\begin{enumerate}[label=(\Alph*)]
    \item $0.692$
    \item $0.800$
    \item $0.900$
    \item $1.000$
\end{enumerate}

\item Compare the $\text{RQ}$ of **palmitic acid** ($\text{C}_{16}\text{H}_{32}\text{O}_2$) with that of **tripalmitin** ($\text{C}_{51}\text{H}_{98}\text{O}_6$). Which statement is mathematically accurate?
\begin{enumerate}[label=(\Alph*)]
    \item Palmitic acid has $\text{RQ} = 16 / 23 = 0.696$, whereas tripalmitin has $\text{RQ} = 102 / 145 = 0.703$; both round to $\approx 0.7$.
    \item Palmitic acid has $\text{RQ} = 1.0$, whereas tripalmitin has $\text{RQ} = 0.5$.
    \item Both have an exact $\text{RQ}$ of $0.900$.
    \item Palmitic acid has $\text{RQ} > 1$, whereas tripalmitin has $\text{RQ} < 1$.
\end{enumerate}

\item In a clinical indirect calorimetry test on a resting human patient, the rate of $\text{O}_2$ consumption is measured as $250\text{ mL/min}$ and the rate of $\text{CO}_2$ production is $200\text{ mL/min}$. What is the patient's Respiratory Quotient ($\text{RQ}$) and what does it indicate about substrate utilization?
\begin{enumerate}[label=(\Alph*)]
    \item $\text{RQ} = 1.25$; pure carbohydrate oxidation
    \item $\text{RQ} = 0.80$; a mixed diet of carbohydrates, fats, and proteins
    \item $\text{RQ} = 0.70$; pure fat oxidation
    \item $\text{RQ} = 0.50$; severe ketosis
\end{enumerate}

\item \textbf{Statement I:} For carbohydrates, the $\text{RQ}$ is $1$ because equal amounts of $\text{CO}_2$ and $\text{O}_2$ are evolved and consumed, respectively. \\
\textbf{Statement II:} Pure proteins or fats are used exclusively as respiratory substrates in most healthy plant tissues.
\begin{enumerate}[label=(\Alph*)]
    \item Both Statement I and Statement II are correct.
    \item Both Statement I and Statement II are incorrect.
    \item Statement I is correct, but Statement II is incorrect.
    \item Statement I is incorrect, but Statement II is correct.
\end{enumerate}

\item Why do **organic acids** (such as malic acid, citric acid, oxalic acid, and tartaric acid) invariably exhibit a Respiratory Quotient ($\text{RQ}$) **greater than $1$**?
\begin{enumerate}[label=(\Alph*)]
    \item Because they contain no carbon atoms.
    \item Because they are already partially oxidized (having a high oxygen-to-carbon ratio compared to carbohydrates), so they require less external $\text{O}_2$ to be completely oxidized to $\text{CO}_2$ and $\text{H}_2\text{O}$.
    \item Because they are oxidized exclusively via anaerobic fermentation.
    \item Because their oxidation consumes $\text{CO}_2$ from the air.
\end{enumerate}

\item Arrange the following respiratory substrates in **descending order** of their Respiratory Quotient ($\text{RQ}$) values:
\begin{enumerate}[label=(\roman*)]
    \item Oxalic acid
    \item Malic acid
    \item Glucose
    \item Albumin (Protein)
    \item Tripalmitin (Fat)
\end{enumerate}
\begin{enumerate}[label=(\Alph*)]
    \item (i) $>$ (ii) $>$ (iii) $>$ (iv) $>$ (v)
    \item (v) $>$ (iv) $>$ (iii) $>$ (ii) $>$ (i)
    \item (iii) $>$ (ii) $>$ (i) $>$ (iv) $>$ (v)
    \item (i) $>$ (iii) $>$ (ii) $>$ (v) $>$ (iv)
\end{enumerate}

\item What is the Respiratory Quotient ($\text{RQ}$) of **succinic acid** ($\text{C}_4\text{H}_6\text{O}_4$) upon complete aerobic respiration?
\begin{enumerate}[label=(\Alph*)]
    \item $0.875$
    \item $1.000$
    \item $1.143$
    \item $1.333$
\end{enumerate}

\item If a plant tissue is respiring **glucose** such that $80\%$ of the glucose is oxidized completely via aerobic respiration and $20\%$ undergoes alcoholic fermentation, what will be the overall measured $\text{RQ}$ of the tissue?
\begin{enumerate}[label=(\Alph*)]
    \item $1.00$
    \item $1.083$
    \item $0.80$
    \item $1.25$
\end{enumerate}

\item Why is the Respiratory Quotient ($\text{RQ}$) of **leaves kept in the dark** typically close to **$1.0$**, whereas illuminated green leaves cannot have their true respiratory $\text{RQ}$ measured directly by simple gas exchange?
\begin{enumerate}[label=(\Alph*)]
    \item Because in the dark, leaves respire stored starch/sucrose ($\text{RQ} = 1$), whereas in the light, simultaneous photosynthetic $\text{CO}_2$ fixation and $\text{O}_2$ evolution mask respiratory gas exchange.
    \item Because leaves stop respiring completely in the presence of light.
    \item Because chlorophyll converts $\text{CO}_2$ into tripalmitin in the dark.
    \item Because stomata close permanently in the dark in all $\text{C}_3$ plants.
\end{enumerate}

\item Consider the complete aerobic oxidation of **pyruvic acid** ($\text{C}_3\text{H}_4\text{O}_3$) in isolated mitochondria. What is the exact Respiratory Quotient ($\text{RQ}$) of pyruvic acid?
\begin{enumerate}[label=(\Alph*)]
    \item $1.00$
    \item $1.20$
    \item $1.33$
    \item $0.80$
\end{enumerate}

\item What is the Respiratory Quotient ($\text{RQ}$) of **acetic acid** ($\text{CH}_3\text{COOH}$ or $\text{C}_2\text{H}_4\text{O}_2$) during complete aerobic oxidation?
\begin{enumerate}[label=(\Alph*)]
    \item $0.70$
    \item $0.90$
    \item $1.00$
    \item $1.50$
\end{enumerate}

\item In **ketogenic diets** or untreated **diabetes mellitus**, the body relies heavily on the oxidation of fatty acids and ketone bodies (such as $\beta$-hydroxybutyrate and acetoacetate). How does this metabolic shift affect the patient's Respiratory Quotient ($\text{RQ}$)?
\begin{enumerate}[label=(\Alph*)]
    \item $\text{RQ}$ rises above $1.0$
    \item $\text{RQ}$ drops toward $0.7$
    \item $\text{RQ}$ becomes infinite ($\infty$)
    \item $\text{RQ}$ remains fixed at $1.0$
\end{enumerate}

\item Calculate the Respiratory Quotient ($\text{RQ}$) of **glycerol** ($\text{C}_3\text{H}_8\text{O}_3$) upon complete aerobic respiration:
\begin{enumerate}[label=(\Alph*)]
    \item $0.700$
    \item $0.857$
    \item $1.000$
    \item $1.167$
\end{enumerate}

\item \textbf{Assertion (A):} In succulent plants like *Bryophyllum*, the leaves taste sour (acidic) in the early morning and become less acidic by late afternoon. \\
\textbf{Reason (R):} At night, incomplete oxidation of carbohydrates accumulates malic acid ($\text{RQ} = 0$), which is subsequently decarboxylated during the day ($\text{RQ} > 1$) to supply $\text{CO}_2$ for the Calvin cycle.
\begin{enumerate}[label=(\Alph*)]
    \item Both (A) and (R) are true and (R) is the correct explanation of (A).
    \item Both (A) and (R) are true but (R) is not the correct explanation of (A).
    \item (A) is true but (R) is false.
    \item Both (A) and (R) are false.
\end{enumerate}

\item What is the Respiratory Quotient ($\text{RQ}$) of **ethanol** ($\text{C}_2\text{H}_5\text{OH}$) when oxidized aerobically by *Acetobacter* or mammalian liver cells to $\text{CO}_2$ and $\text{H}_2\text{O}$?
\begin{enumerate}[label=(\Alph*)]
    \item $0.667$
    \item $0.750$
    \item $1.000$
    \item $1.333$
\end{enumerate}

\item During **glycolysis** alone (conversion of $1\text{ molecule}$ of glucose to $2\text{ molecules}$ of pyruvic acid), why can a Respiratory Quotient ($\text{RQ}$) not be calculated?
\begin{enumerate}[label=(\Alph*)]
    \item Because glycolysis neither consumes $\text{O}_2$ nor evolves $\text{CO}_2$ ($0 / 0$).
    \item Because glycolysis consumes $6\text{ O}_2$ without evolving $\text{CO}_2$.
    \item Because glycolysis evolves $6\text{ CO}_2$ without consuming $\text{O}_2$.
    \item Because glycolysis produces tripalmitin.
\end{enumerate}

\item Which of the following pairs of substrates have the **exact same** theoretical Respiratory Quotient ($\text{RQ}$)?
\begin{enumerate}[label=(\Alph*)]
    \item Glucose ($\text{C}_6\text{H}_{12}\text{O}_6$) and Lactic acid ($\text{C}_3\text{H}_6\text{O}_3$)
    \item Tripalmitin ($\text{C}_{51}\text{H}_{98}\text{O}_6$) and Malic acid ($\text{C}_4\text{H}_6\text{O}_5$)
    \item Oxalic acid ($\text{C}_2\text{H}_2\text{O}_4$) and Citric acid ($\text{C}_6\text{H}_8\text{O}_7$)
    \item Palmitic acid ($\text{C}_{16}\text{H}_{32}\text{O}_2$) and Sucrose ($\text{C}_{12}\text{H}_{22}\text{O}_{11}$)
\end{enumerate}

\item A student places $10\text{ g}$ of germinating **groundnut** seeds in Respirometer A and $10\text{ g}$ of germinating **maize** grains in Respirometer B under identical aerobic conditions. After $1\text{ hour}$, which observation is correct?
\begin{enumerate}[label=(\Alph*)]
    \item Respirometer A shows $\text{RQ} \approx 1.0$, while Respirometer B shows $\text{RQ} \approx 0.7$.
    \item Respirometer A shows $\text{RQ} \approx 0.7$ (fat-rich seeds), while Respirometer B shows $\text{RQ} \approx 1.0$ (starch-rich grains).
    \item Both respirometers show $\text{RQ} = 0.0$.
    \item Both respirometers show $\text{RQ} = 1.33$.
\end{enumerate}

\end{enumerate}

\newpage

\section*{Answer Key}

\begin{multicols}{5}
\begin{enumerate}[label=\textbf{\arabic*.}]
    \item (B)
    \item (C)
    \item (A)
    \item (A)
    \item (D)
    \item (C)
    \item (A)
    \item (B)
    \item (C)
    \item (D)
    \item (A)
    \item (A)
    \item (A)
    \item (B)
    \item (C)
    \item (B)
    \item (B)
    \item (B)
    \item (B)
    \item (A)
    \item (C)
    \item (C)
    \item (B)
    \item (B)
    \item (B)
    \item (A)
    \item (A)
    \item (A)
    \item (A)
    \item (B)
    \item (C)
    \item (B)
    \item (A)
    \item (C)
    \item (B)
    \item (A)
    \item (B)
    \item (C)
    \item (B)
    \item (B)
    \item (A)
    \item (A)
    \item (A)
    \item (A)
    \item (B)
\end{enumerate}
\end{multicols}

\section*{Detailed Solutions}

\begin{enumerate}[label=\textbf{\arabic*.}]

\item \textbf{(B)} Verbatim from NCERT Class 11 Biology (Respiration in Plants):
$$\text{RQ} = \frac{\text{Volume of } \text{CO}_2 \text{ evolved}}{\text{Volume of } \text{O}_2 \text{ consumed}}$$

\item \textbf{(C)} According to NCERT, when fats are used in respiration, the $\text{RQ}$ is less than $1$. Calculations for a fatty acid, **tripalmitin**, give an $\text{RQ}$ of **$0.7$** ($102 / 145 = 0.703 \approx 0.7$).

\item \textbf{(A)} Verbatim equation from NCERT Class 11 Biology:
$$2(\text{C}_{51}\text{H}_{98}\text{O}_6) + 145\text{O}_2 \longrightarrow 102\text{CO}_2 + 98\text{H}_2\text{O} + \text{Energy}$$
Hence $x = 145$, $y = 102$, and $z = 98$.

\item \textbf{(A)} Fats are highly reduced hydrocarbons with very little oxygen in their molecular structure (e.g., only $6\text{ O}$ atoms for $51\text{ C}$ and $98\text{ H}$ in tripalmitin). Consequently, a large amount of external $\text{O}_2$ is required to oxidize all the carbon and hydrogen, making the denominator ($\text{O}_2$ consumed) much larger than the numerator ($\text{CO}_2$ evolved), yielding $\text{RQ} \approx 0.7$.

\item \textbf{(D)} Balanced equation for malic acid ($\text{C}_4\text{H}_6\text{O}_5$):
$$\text{C}_4\text{H}_6\text{O}_5 + 3\text{O}_2 \longrightarrow 4\text{CO}_2 + 3\text{H}_2\text{O}$$
$$\text{RQ} = \frac{4\text{ CO}_2}{3\text{ O}_2} = \mathbf{1.33}$$

\item \textbf{(C)} Balanced equation for oxalic acid ($\text{C}_2\text{H}_2\text{O}_4$):
$$2\text{C}_2\text{H}_2\text{O}_4 + \text{O}_2 \longrightarrow 4\text{CO}_2 + 2\text{H}_2\text{O} \quad \left(\text{or } \text{C}_2\text{H}_2\text{O}_4 + 0.5\text{O}_2 \longrightarrow 2\text{CO}_2 + \text{H}_2\text{O}\right)$$
$$\text{RQ} = \frac{4\text{ CO}_2}{1\text{ O}_2} = \mathbf{4.00}$$

\item \textbf{(A)} Starchy wheat seeds respire carbohydrates ($\text{RQ} = 1.0$, P-3). Fatty castor/groundnut seeds respire oils ($\text{RQ} = 0.7$, Q-1). Proteinaceous pea/bean seeds have $\text{RQ} \approx 0.9$ (R-2). Succulent CAM plants like *Opuntia* at night fix all $\text{CO}_2$ into malic acid without releasing $\text{CO}_2$ ($\text{RQ} = 0.0$, S-4).

\item \textbf{(B)} At night, CAM plants (*Opuntia*) open their stomata, absorb $\text{O}_2$ to partially oxidize starch into PEP and malic acid, and simultaneously fix all respiratory and atmospheric $\text{CO}_2$ via PEP carboxylase. Because $0\text{ CO}_2$ is evolved while $\text{O}_2$ is consumed, $\text{RQ} = 0 / \text{O}_2 = \mathbf{0}$.

\item \textbf{(C)} During the day, *Opuntia* closes its stomata and decarboxylates the stored malic acid ($\text{C}_4\text{H}_6\text{O}_5$), releasing large amounts of $\text{CO}_2$ internally while consuming relatively little external $\text{O}_2$, resulting in **$\text{RQ} > 1$** ($\approx 1.33$).

\item \textbf{(D)} In alcoholic fermentation ($\text{C}_6\text{H}_{12}\text{O}_6 \rightarrow 2\text{C}_2\text{H}_5\text{OH} + 2\text{CO}_2$), $2\text{ moles}$ of $\text{CO}_2$ are evolved while **$0\text{ moles}$ of $\text{O}_2$** are consumed. Therefore:
$$\text{RQ} = \frac{2}{0} = \mathbf{\infty \text{ (Infinity)}}$$

\item \textbf{(A)} In pure lactic acid fermentation ($\text{C}_6\text{H}_{12}\text{O}_6 \rightarrow 2\text{CH}_3\text{CH(OH)COOH}$), neither $\text{O}_2$ is consumed ($0$) nor $\text{CO}_2$ is evolved ($0$), giving $0/0$ (undefined).

\item \textbf{(A)} Both statements are verbatim from NCERT Class 11 Biology (Section 14.5: *Respiratory Quotient*): \textit{``The respiratory quotient depends upon the type of respiratory substrate used during respiration... It is also important to recognise that in living organisms respiratory substrates are often more than one; pure proteins or fats are never used as respiratory substrates.''}

\item \textbf{(A)} Balanced equation for oleic acid ($\text{C}_{18}\text{H}_{34}\text{O}_2$):
$$\text{C}_{18}\text{H}_{34}\text{O}_2 + 25.5\text{O}_2 \longrightarrow 18\text{CO}_2 + 17\text{H}_2\text{O}$$
$$\text{RQ} = \frac{18}{25.5} = \frac{36}{51} = \mathbf{0.706}$$

\item \textbf{(B)} During peak germination of oilseeds, fatty acids ($\text{O}$-poor) are converted into sucrose ($\text{O}$-rich) via the glyoxylate cycle and gluconeogenesis. Incorporating oxygen atoms into newly formed sugar molecules requires significant external $\text{O}_2$ uptake without releasing $\text{CO}_2$, depressing the $\text{RQ}$ to $\sim 0.4\text{--}0.5$.

\item \textbf{(C)} Calculating $\text{RQ}$:
$$\text{RQ} = \frac{\text{Volume of } \text{CO}_2 \text{ evolved}}{\text{Volume of } \text{O}_2 \text{ consumed}} = \frac{84\text{ mL}}{120\text{ mL}} = \mathbf{0.7}$$
An $\text{RQ}$ of $0.7$ indicates that **fats/oils** are the primary respiratory substrate.

\item \textbf{(B)} For wheat seeds ($\text{RQ} = 1$), the volume of $\text{CO}_2$ evolved equals the volume of $\text{O}_2$ absorbed. Without $\text{KOH}$ to absorb $\text{CO}_2$, total gas volume inside the respirometer remains unchanged, so there is **no net movement** of the liquid level.

\item \textbf{(B)} For castor seeds ($\text{RQ} = 0.7$), the volume of $\text{O}_2$ absorbed ($10\text{ units}$) exceeds the volume of $\text{CO}_2$ evolved ($7\text{ units}$). Without $\text{KOH}$, total gas volume decreases, creating a partial vacuum that causes the liquid level to **rise**.

\item \textbf{(B)} When $\text{RQ} > 1$ (organic acids), the volume of $\text{CO}_2$ evolved exceeds the volume of $\text{O}_2$ absorbed. Without $\text{KOH}$, total gas pressure increases inside the chamber, pushing the liquid column **downward (fall)**.

\item \textbf{(B)} Proteins have an intermediate degree of reduction (more oxygen than fats, less than carbohydrates), and during amino acid catabolism, amino groups are removed and converted into urea/uric acid (which consumes $\text{CO}_2$ in the urea cycle), giving a net $\text{RQ}$ of **$\approx 0.9$** ($0.8\text{--}0.9$).

\item \textbf{(A)} For complete combustion of $\text{C}_x\text{H}_y\text{O}_z$:
$$\text{C}_x\text{H}_y\text{O}_z + \left(x + \frac{y}{4} - \frac{z}{2}\right)\text{O}_2 \longrightarrow x\text{CO}_2 + \frac{y}{2}\text{H}_2\text{O}$$
Therefore, $\text{RQ} = \dfrac{x}{x + \frac{y}{4} - \frac{z}{2}}$.

\item \textbf{(C)} For citric acid ($\text{C}_6\text{H}_8\text{O}_7$):
$$\text{O}_2 \text{ required} = 6 + \frac{8}{4} - \frac{7}{2} = 6 + 2 - 3.5 = 4.5\text{ O}_2$$
$$\text{RQ} = \frac{6\text{ CO}_2}{4.5\text{ O}_2} = \frac{12}{9} = \mathbf{1.33}$$

\item \textbf{(C)} For tartaric acid ($\text{C}_4\text{H}_6\text{O}_6$):
$$\text{O}_2 \text{ required} = 4 + \frac{6}{4} - \frac{6}{2} = 4 + 1.5 - 3 = 2.5\text{ O}_2$$
$$\text{RQ} = \frac{4\text{ CO}_2}{2.5\text{ O}_2} = \frac{8}{5} = \mathbf{1.60}$$

\item \textbf{(B)} When $50\%$ of $\text{O}_2$ is used to oxidize carbohydrate ($\text{RQ} = 1.0$) and $50\%$ is used to oxidize fat ($\text{RQ} = 0.7$):
$$\text{Mixed RQ} = (0.5 \times 1.0) + (0.5 \times 0.7) = 0.50 + 0.35 = \mathbf{0.85}$$

\item \textbf{(B)} In the well-fed state, carbohydrates are oxidized ($\text{RQ} = 1.0$). During fasting/starvation, glycogen is depleted within $24\text{ hours}$ and the body shifts to adipose triglycerides ($\text{RQ} = 0.7$). Once fat stores are exhausted in terminal starvation, structural proteins are catabolized ($\text{RQ} \approx 0.8\text{--}0.9$).

\item \textbf{(B)} Converting oxygen-rich carbohydrates (glucose) into oxygen-poor fatty acids releases excess oxygen atoms internally (used in oxidative pathways) and evolves $\text{CO}_2$ during pyruvate decarboxylation to Acetyl-CoA without requiring external $\text{O}_2$, pushing $\text{RQ}$ **above $1.0$** ($1.1\text{--}1.3$).

\item \textbf{(A)} Before the seed coat (testa) ruptures, water-imbibed tissues restrict $\text{O}_2$ entry. The embryo performs partial anaerobic fermentation (releasing extra $\text{CO}_2$ without $\text{O}_2$ uptake) alongside aerobic respiration, causing $\text{RQ} > 1$. Once the testa splits and $\text{O}_2$ enters freely, $\text{RQ}$ drops to $1.0$.

\item \textbf{(A)} **Sucrose** ($\text{C}_{12}\text{H}_{22}\text{O}_{11}$) is a carbohydrate:
$$\text{C}_{12}\text{H}_{22}\text{O}_{11} + 12\text{O}_2 \longrightarrow 12\text{CO}_2 + 11\text{H}_2\text{O} \implies \text{RQ} = \frac{12}{12} = \mathbf{1.0}$$

\item \textbf{(A)} For stearic acid ($\text{C}_{18}\text{H}_{36}\text{O}_2$):
$$\text{C}_{18}\text{H}_{36}\text{O}_2 + 26\text{O}_2 \longrightarrow 18\text{CO}_2 + 18\text{H}_2\text{O} \implies \text{RQ} = \frac{18}{26} = \mathbf{0.692}$$

\item \textbf{(A)} For palmitic acid ($\text{C}_{16}\text{H}_{32}\text{O}_2$): $\text{C}_{16}\text{H}_{32}\text{O}_2 + 23\text{O}_2 \rightarrow 16\text{CO}_2 + 16\text{H}_2\text{O} \implies \text{RQ} = 16/23 = \mathbf{0.696}$. For tripalmitin ($\text{C}_{51}\text{H}_{98}\text{O}_6$): $\text{RQ} = 102/145 = \mathbf{0.703}$. Both round to $\approx 0.7$.

\item \textbf{(B)} $\text{RQ} = \dfrac{200\text{ mL/min}}{250\text{ mL/min}} = \mathbf{0.80}$, which is the classic physiological $\text{RQ}$ of a person metabolizing a balanced mixed diet of carbohydrates, fats, and proteins.

\item \textbf{(C)} Statement I is correct (verbatim NCERT). Statement II is incorrect because NCERT explicitly states that pure proteins or fats are **never** used as exclusive respiratory substrates in living organisms.

\item \textbf{(B)} Organic acids (e.g., malic acid $\text{C}_4\text{H}_6\text{O}_5$, oxalic acid $\text{C}_2\text{H}_2\text{O}_4$) possess a higher proportion of oxygen atoms per carbon atom than carbohydrates ($\text{CH}_2\text{O}$). Because they are already partially oxidized, they require less external $\text{O}_2$ for complete oxidation, resulting in $\text{RQ} > 1$.

\item \textbf{(A)} Exact values:
\begin{itemize}
    \item Oxalic acid (i): $4.00$
    \item Malic acid (ii): $1.33$
    \item Glucose (iii): $1.00$
    \item Albumin/Protein (iv): $\approx 0.90$
    \item Tripalmitin/Fat (v): $\approx 0.70$
\end{itemize}
Descending order: **(i) $>$ (ii) $>$ (iii) $>$ (iv) $>$ (v)**.

\item \textbf{(C)} For succinic acid ($\text{C}_4\text{H}_6\text{O}_4$):
$$\text{O}_2 \text{ required} = 4 + \frac{6}{4} - \frac{4}{2} = 4 + 1.5 - 2 = 3.5\text{ O}_2$$
$$\text{RQ} = \frac{4\text{ CO}_2}{3.5\text{ O}_2} = \frac{8}{7} = \mathbf{1.143}$$

\item \textbf{(B)} Consider $10\text{ moles}$ of glucose ($8\text{ aerobic} + 2\text{ anaerobic}$):
\begin{itemize}
    \item $8\text{ aerobic glucose}$ consume $8 \times 6 = 48\text{ O}_2$ and evolve $8 \times 6 = 48\text{ CO}_2$.
    \item $2\text{ anaerobic glucose}$ consume $0\text{ O}_2$ and evolve $2 \times 2 = 4\text{ CO}_2$.
\end{itemize}
Total $\text{CO}_2 = 48 + 4 = 52$; Total $\text{O}_2 = 48$.
$$\text{RQ} = \frac{52}{48} = \mathbf{1.083}$$

\item \textbf{(A)} In the dark, photosynthesis ceases and leaves respire stored carbohydrates ($\text{RQ} = 1.0$). In the light, photosynthesis consumes respiratory $\text{CO}_2$ and releases $\text{O}_2$, obscuring true respiratory gas exchange.

\item \textbf{(B)} Balanced equation for pyruvic acid ($\text{C}_3\text{H}_4\text{O}_3$):
$$\text{C}_3\text{H}_4\text{O}_3 + 2.5\text{O}_2 \longrightarrow 3\text{CO}_2 + 2\text{H}_2\text{O} \implies \text{RQ} = \frac{3}{2.5} = \mathbf{1.20}$$

\item \textbf{(C)} Acetic acid ($\text{C}_2\text{H}_4\text{O}_2$) has the empirical formula $\text{CH}_2\text{O}$ (identical to carbohydrates):
$$\text{C}_2\text{H}_4\text{O}_2 + 2\text{O}_2 \longrightarrow 2\text{CO}_2 + 2\text{H}_2\text{O} \implies \text{RQ} = \frac{2}{2} = \mathbf{1.00}$$

\item \textbf{(B)} In untreated diabetes mellitus or ketogenic diets, glucose utilization is impaired/restricted and tissues oxidize fatty acids ($\text{RQ} = 0.7$) and ketone bodies ($\beta$-hydroxybutyrate $\text{RQ} = 0.89$; acetoacetate $\text{RQ} = 1.0$), driving the overall $\text{RQ}$ down toward **$0.7$**.

\item \textbf{(B)} Balanced equation for glycerol ($\text{C}_3\text{H}_8\text{O}_3$):
$$\text{C}_3\text{H}_8\text{O}_3 + 3.5\text{O}_2 \longrightarrow 3\text{CO}_2 + 4\text{H}_2\text{O} \implies \text{RQ} = \frac{3}{3.5} = \frac{6}{7} = \mathbf{0.857}$$

\item \textbf{(A)} In *Bryophyllum* (a CAM succulent), nocturnal acidification accumulates malic acid in the vacuoles ($\text{RQ} = 0$, making leaves taste sour in the morning), while daytime deacidification decarboxylates malic acid ($\text{RQ} > 1$, lowering leaf acidity by afternoon).

\item \textbf{(A)} Balanced equation for ethanol ($\text{C}_2\text{H}_6\text{O}$):
$$\text{C}_2\text{H}_6\text{O} + 3\text{O}_2 \longrightarrow 2\text{CO}_2 + 3\text{H}_2\text{O} \implies \text{RQ} = \frac{2}{3} = \mathbf{0.667}$$

\item \textbf{(A)} Glycolysis ($\text{C}_6\text{H}_{12}\text{O}_6 + 2\text{NAD}^+ + 2\text{ADP} + 2\text{P}_i \rightarrow 2\text{C}_3\text{H}_4\text{O}_3 + 2\text{NADH} + 2\text{H}^+ + 2\text{ATP}$) neither consumes $\text{O}_2$ nor releases $\text{CO}_2$, making $\text{RQ}$ undefined ($0/0$).

\item \textbf{(A)} Both glucose ($\text{C}_6\text{H}_{12}\text{O}_6$) and lactic acid ($\text{C}_3\text{H}_6\text{O}_3$) share the empirical formula $(\text{CH}_2\text{O})_n$. For complete aerobic oxidation: $\text{C}_3\text{H}_6\text{O}_3 + 3\text{O}_2 \rightarrow 3\text{CO}_2 + 3\text{H}_2\text{O}$, giving $\text{RQ} = 3/3 = \mathbf{1.0}$, identical to glucose ($6/6 = 1.0$).

\item \textbf{(B)} Groundnut seeds are rich in stored fats/oils ($\text{RQ} \approx 0.7$), whereas maize grains (caryopses) have a large starchy endosperm ($\text{RQ} \approx 1.0$).

\end{enumerate}

\end{document}
</USER_REQUEST>
<ADDITIONAL_METADATA>
The current local time is: 2026-11-16T12:04:27+05:30.

The user's current state is as follows:
Active Document: /Users/laxmikumari/Desktop/web/testseries/scripts/add_botany_zoology_to_part_11.mjs (LANGUAGE_UNSPECIFIED)
Cursor is on line: 1
Other open documents:
- /Users/laxmikumari/Desktop/web/testseries/.env.local (LANGUAGE_UNSPECIFIED)
- /Users/laxmikumari/Downloads/neet_five_kingdom_classification_test.tex (LANGUAGE_UNSPECIFIED)
- /Users/laxmikumari/Downloads/GOC_Extracted_Questions.tex (LANGUAGE_UNSPECIFIED)
- /Users/laxmikumari/Downloads/neet_biology_question_bank_phases_of_growth_set_2.tex (LANGUAGE_UNSPECIFIED)
- /Users/laxmikumari/Desktop/web/testseries/src/app/api/admin/notifications/route.js (LANGUAGE_UNSPECIFIED)
</ADDITIONAL_METADATA>